\chapter{Fundamente teoretice și tehnologiile utilizate}

Acest capitol prezintă baza teoretică și tehnologică pe care se sprijină aplicația WhereWeFishin. Dacă în capitolul introductiv am descris motivația și obiectivele proiectului, acum mă concentrez pe conceptele care au ghidat proiectarea și implementarea efectivă a sistemului. Sunt analizate arhitectura generală, modelul de date, componenta geospațială, fluxul de rezervare, autentificarea și componenta de analiză media bazată pe inteligență artificială.

Scopul nu este reproducerea detaliilor de cod, ci clarificarea principiilor pe care funcționează sistemul și justificarea alegerilor tehnologice. Această abordare urmează ideea clasică din ingineria software că structura unui sistem trebuie dedusă din cerințe și din separarea clară a responsabilităților \citep{sommerville2015,pressman2014}.

\section{Arhitectura generală a sistemului}

O aplicație care deservește simultan utilizatori finali, angajați și administratori are nevoie de o arhitectură bine structurată. Trebuie să gestioneze interacțiuni de interfață, validări de business, persistența datelor, comunicarea cu servicii externe și, în cazul de față, procesarea unor fișiere media. Din acest motiv, WhereWeFishin este construită ca un sistem cu mai multe componente distincte, dar integrate coerent.

Interfața cu utilizatorul este o aplicație web de tip SPA dezvoltată cu Angular \citep{angular}. Această alegere este potrivită pentru o aplicație cu navigare fluidă, componente reutilizabile, formulare complexe și actualizare dinamică a conținutului. WhereWeFishin are fluxuri variate — autentificare, hartă interactivă, rezervare, administrare, upload de fișiere — și o interfață bazată pe componente permite separarea clară a responsabilităților.

Backend-ul este un API HTTP realizat cu ASP.NET Core \citep{aspnetcore}. El centralizează logica de business, aplică regulile de validare, controlează accesul la resurse și expune servicii pentru interfața web. ASP.NET Core oferă suport solid pentru autentificare, autorizare, middleware și interoperabilitate cu servicii externe. Stilul arhitectural adoptat este REST, larg utilizat în aplicațiile web moderne \citep{fielding2000}.

Persistența datelor se face printr-o bază de date relațională, accesată prin Entity Framework Core \citep{efcore,sqlserver}. Aceste instrumente mapează entitățile domeniului la tabele și permit interogări clare și stabile. Alegerea unui model relațional este justificată de numărul mare de relații dintre date: utilizatori, locuri de pescuit, pontoane, sesiuni de pescuit, recenzii și rezultate ale analizelor media.

Un element distinctiv al arhitecturii este microserviciul separat pentru analiza imaginilor și videoclipurilor, implementat în Python cu Flask \citep{flask}. Separarea sa de API-ul principal are o justificare practică: procesarea media implică biblioteci specializate, cerințe diferite de execuție și timpi mai lungi de răspuns. Prin izolarea acestei responsabilități, se reduce riscul ca o operație costisitoare să afecteze stabilitatea backend-ului principal.

Modul în care componentele comunică este simplu: utilizatorul interacționează cu interfața Angular, aceasta trimite cereri către API-ul .NET, iar API-ul aplică regulile de business și salvează datele în baza de date. Când un fișier media trebuie analizat, backend-ul îl trimite microserviciului Python, primește rezultatele și le stochează.

\section{Modelul de date și entitățile de domeniu}

Un sistem informatic orientat spre un domeniu concret trebuie să pornească de la o modelare corectă a obiectelor care definesc acea activitate. În cazul unei platforme pentru pescuit recreativ, modelul de date nu descrie doar utilizatori și informații generale, ci trebuie să reflecte realități operaționale: locuri de pescuit cu coordonate geografice, zone rezervabile, sesiuni planificate și fluxuri de validare.

Entitatea centrală este \textbf{locul de pescuit}. El conține informații descriptive (denumire, descriere, imagine), coordonate geografice, tarif orar și legături cu utilizatorul care l-a creat sau cu managerul care îl administrează. În jurul lui gravitează celelalte entități importante: pontoanele, recenziile, sesiunile de pescuit, angajații și repopulările cu pești.

\textbf{Pontonul} este o subzonă rezervabilă din interiorul unei locații. Are o dimensiune spațială precisă, descrisă prin coordonate geografice, ceea ce permite reprezentarea lui pe hartă și rezervarea explicită. Această alegere simplifică administrarea spațiului și oferă utilizatorului o imagine mai clară a ceea ce rezervă.

\textbf{Sesiunea de pescuit} este entitatea care materializează o rezervare. Ea leagă utilizatorul de o locație, de un ponton (dacă este cazul), de un interval temporal, de un cost și de o stare operațională. Termenul de „sesiune de pescuit" este mai precis decât „rezervare" pentru descrierea obiectului intern, deoarece include toată informația necesară gestionării operaționale a procesului.

\textbf{Utilizatorul} este prezent în aproape toate fluxurile platformei: inițiază rezervări, lasă recenzii, poate deveni manager. \textbf{Recenziile} oferă feedback despre locații. \textbf{Repopulările cu pești} descriu acțiuni administrative specifice domeniului. \textbf{Analizele media} stochează rezultatele generate de microserviciul Python.

Figura~\ref{fig:model-date} prezintă principalele entități și relațiile dintre ele.

\begin{figure}[ht]
\centering
\begin{tikzpicture}[
  every node/.style={font=\small},
  ent/.style={draw, rounded corners, align=center, minimum width=2.9cm, minimum height=0.85cm, fill=gray!10},
  arrow/.style={-Latex, thick},
  lbl/.style={font=\footnotesize, fill=white, inner sep=1pt}
]
\node[ent] (user)    at (0, 0)    {Utilizator};
\node[ent] (spot)    at (5, 0)    {Loc de pescuit};
\node[ent] (pontoon) at (5, -2.2) {Ponton};
\node[ent] (session) at (0, -2.2) {Sesiune de pescuit};
\node[ent] (review)  at (5, -4.4) {Recenzie};
\node[ent] (analysis)at (0, -4.4) {Analiză media};
\node[ent] (restock) at (2.5,-6.2){Repopulare pești};

\draw[arrow] (user)    -- node[lbl, above]{face} (session);
\draw[arrow] (session) -- node[lbl, above]{la} (spot);
\draw[arrow] (spot)    -- node[lbl, right]{are} (pontoon);
\draw[arrow] (session) -- node[lbl, right]{include} (pontoon);
\draw[arrow] (user)    -- node[lbl, left]{scrie} (review);
\draw[arrow] (review)  -- node[lbl, right]{despre} (spot);
\draw[arrow] (user)    -- node[lbl, left]{încarcă} (analysis);
\draw[arrow] (spot)    -- node[lbl, right]{primește} (restock);
\end{tikzpicture}
\caption{Principalele entități și relații din modelul de date}
\label{fig:model-date}
\end{figure}

Privit în ansamblu, modelul de date răspunde nevoii de coerență între lumea reală și reprezentarea sa informatică. Utilizarea unei baze de date relaționale și a unui ORM este potrivită deoarece permite exprimarea clară a relațiilor și reduce riscul inconsistențelor.

\section{Harta interactivă și componenta geospațială}

În domeniul pescuitului recreativ, poziționarea geografică nu este un detaliu secundar. Utilizatorul nu caută doar un nume de locație, ci vrea să vadă unde se află aceasta, ce locuri există în apropiere și cum este organizat spațiul.

Aplicația folosește biblioteca Leaflet pentru harta interactivă \citep{leaflet}. Aceasta este ușoară, bine documentată, compatibilă cu dispozitivele mobile și suportă marcarea locațiilor, desenarea zonelor, ferestre informative și actualizare dinamică a conținutului. Locațiile sunt descrise prin coordonate de latitudine și longitudine, suficiente pentru afișarea pe hartă și pentru operații de tip centrare sau zoom.

Harta funcționează în două scenarii principale. Primul este explorarea: utilizatorul caută un loc de pescuit, vede pozițiile disponibile pe hartă și decide ce vrea să rezerve. Al doilea este administrarea: managerul delimitează sau actualizează vizual zona unui ponton, controlând cum apare aceasta în interfața publică.

Un aspect important este că pontonul nu este tratat doar ca o etichetă, ci ca o zonă cu poziționare geografică explicită. Utilizatorul poate înțelege mai bine ce rezervă, iar managerul poate controla mai precis împărțirea spațiului. Într-un domeniu în care proximitatea față de apă sau organizarea pe mal influențează decizia, o astfel de reprezentare este utilă.

\section{Rezervări și sesiuni de pescuit}

Fluxul de rezervare transformă informația statică despre o locație într-o interacțiune concretă. Utilizatorul solicită utilizarea unui loc sau a unui ponton pentru un interval de timp determinat. Rezultatul intern al acestei acțiuni este o sesiune de pescuit — entitatea care reține toate datele necesare gestionării operaționale.

Această distincție între „rezervare" (acțiunea utilizatorului) și „sesiune de pescuit" (obiectul intern al sistemului) este utilă atât tehnic, cât și pentru claritatea textului. Sistemul urmărește un obiect bine definit care leagă utilizatorul de locație, ponton, interval temporal, cost și stare operațională.

Pentru crearea corectă a unei sesiuni, sistemul verifică mai multe condiții: existența locației, apartenența pontonului la locația selectată și absența conflictelor temporale. Tratarea corectă a datei și orei este importantă mai ales când utilizatorii interacționează de pe dispozitive diferite.

Costul este calculat pe baza duratei și a tarifului orar al locației. Includerea lui în entitatea sesiunii permite păstrarea istoricului și analiza rezervărilor independent de valoarea curentă a tarifului.

Când plata online este activată, rezervarea se completează prin integrarea cu Stripe \citep{stripe}. Aplicația nu implementează logica de plată de la zero, ci delegă această responsabilitate unui serviciu specializat, păstrând controlul asupra validărilor de business.

După confirmare, sesiunea intră într-un flux de validare la fața locului prin cod QR și token de verificare. Această soluție reduce dependența de confirmări verbale și creează un proces controlat, ușor de urmărit de angajați și manageri.

\section{Autentificare, autorizare și securitate}

Orice aplicație care gestionează conturi, date personale și rezervări trebuie să trateze securitatea ca o cerință de bază. La WhereWeFishin, această nevoie este accentuată de faptul că sistemul deservește categorii diferite de utilizatori, fiecare cu permisiuni distincte.

Autentificarea se bazează pe tokenuri JWT, standardizate prin RFC 7519 \citep{jwt}. Acest mecanism este potrivit pentru un API consumat de o interfață SPA: permite transportul controlat al informațiilor despre utilizator și susține un model fără stare pe server, compatibil cu serviciile REST.

Autentificarea singură nu este suficientă. Sistemul trebuie să decidă și ce poate face fiecare utilizator autentificat — aceasta este autorizarea. Un utilizator obișnuit poate face rezervări și consulta propriul istoric, dar nu poate administra locații ale altor persoane. Un manager operează numai asupra resurselor asociate locațiilor sale. ASP.NET Core oferă suport explicit pentru aceste mecanisme, inclusiv limitarea accesului la rute pe baza rolurilor \citep{aspnetcore}.

Securitatea nu se reduce la un singur mecanism. Ea include validarea datelor primite, protejarea fluxurilor sensibile și delimitarea clară a interacțiunilor dintre componente. De exemplu, operațiile privind plățile sau accesul la rezultatele analizelor media sunt expuse doar în condiții bine definite.

\section{Analiza imaginilor și videoclipurilor cu ajutorul inteligenței artificiale}

Una dintre componentele care diferențiază WhereWeFishin de o platformă obișnuită de rezervări este integrarea analizei media bazate pe inteligență artificială. Această zonă aparține domeniului viziunii artificiale — ramura AI care extrage automat informații din imagini și secvențe video. Interesul principal este detectarea și identificarea speciilor de pești în fișierele încărcate de utilizator.

În analiza imaginilor statice, problema de bază este detecția obiectelor: localizarea peștelui în imagine și clasificarea sa. Familia de modele YOLO este larg utilizată în acest domeniu datorită vitezei și eficienței practice \citep{ultralytics,redmon2016yolo}. Pentru WhereWeFishin, aceasta este alegerea potrivită deoarece funcționează bine atât pe imagini, cât și pe secvențe video.

Rezultatul unei detecții nu este o simplă decizie de tip „există" sau „nu există" pește. Modelul returnează poziția în imagine, gradul de încredere al clasificării și specia estimată. Aceste informații permit atât o reprezentare vizuală a rezultatului, cât și generarea unor statistici relevante.

Pentru videoclipuri, analiza devine mai complexă. Dacă fiecare cadru ar fi analizat independent, același pește ar putea fi contorizat de mai multe ori. De aceea, în WhereWeFishin este folosit mecanismul ByteTrack pentru urmărirea obiectelor în timp \citep{zhang2022bytetrack}. Acesta menține o asociere coerentă între detecțiile succesive și reduce fragmentarea traseelor când același exemplar apare în mai multe cadre.

Procesarea video implică mai mulți pași: citirea cadrelor, analiza fiecărui cadru, corelarea temporală a detecțiilor și stocarea rezultatelor. Bibliotecile OpenCV și FFmpeg sunt folosite pentru prelucrarea cadrelor și recodarea fișierelor video \citep{opencv,ffmpeg}.

Rezultatul final poate include mai multe categorii de date: numărul total de detecții, specia dominantă, distribuția observațiilor și momentele în care au apărut obiectele detectate. Această bogăție de informație permite atât o prezentare intuitivă în interfață, cât și o interpretare mai detaliată.

Este important de menționat că analiza automată nu este infailibilă. Performanța depinde de calitatea imaginilor, iluminare, claritatea apei, unghiuri de filmare și gradul de adaptare al modelului la speciile urmărite. Rezultatele trebuie înțelese în raport cu limitele normale ale sistemelor de viziune artificială.

\section{Servicii auxiliare și integrarea cu sisteme externe}

Pe lângă componentele principale, platforma are nevoie și de servicii auxiliare care completează experiența utilizatorului și susțin operațiile curente.

Integrarea cu Stripe pentru plăți online reduce complexitatea internă și delegă logica sensibilă a tranzacțiilor unui serviciu specializat \citep{stripe}. Aceasta permite și extinderea ulterioară a fluxului de plată fără modificarea arhitecturii principale.

Sistemul de notificări prin e-mail trimite confirmări după rezervare sau mesaje legate de operații importante. Notificările cresc claritatea interacțiunii și reduc dependența utilizatorului de a fi conectat permanent la aplicație pentru a afla rezultatul unei acțiuni.

Aceste servicii auxiliare au și un rol de integrare între module. O rezervare confirmată poate genera simultan o înregistrare internă, o notificare externă și o referință pentru validarea prin cod QR. O analiză media finalizată trebuie să fie accesibilă în interfață printr-un URL stabil. Aceste detalii nu schimbă modelul de domeniu, dar influențează direct calitatea soluției finale.

\section{Sinteza capitolului}

Fundamentele prezentate în acest capitol conturează baza pe care funcționează WhereWeFishin. Arhitectura separată pe componente, modelul de date orientat spre entitățile reale ale domeniului, harta interactivă, fluxurile de rezervare și validare, mecanismele de securitate și componenta de analiză media nu sunt elemente independente, ci părți ale aceluiași sistem.

Înțelegerea acestor fundamente este importantă pentru capitolele următoare, deoarece explică de ce anumite decizii de proiectare și implementare au fost luate și cum contribuie ele la obținerea unei platforme coerente.
